\documentclass[11pt,a4paper]{article}
\usepackage[margin=2.5cm]{geometry}
\usepackage[utf8]{inputenc}
\usepackage[T1]{fontenc}
\usepackage[french]{babel}
\usepackage{hyperref}

\begin{document}
	\pagestyle{empty}
	
	\noindent
	Friedly WOLI \\
	Toulouse, 31100 \\
	0767339290 \\
	friedlysedjro@gmail.com
	
	\begin{flushright}
		Capgemini Engineering \\
		Toulouse / Aix-en-Provence
	\end{flushright}
	
	\bigskip
	
	\textbf{Objet : Candidature – Alternance en calcul mécanique} \\
	
	\bigskip
	
	Madame, Monsieur, \\
	
	Actuellement étudiant en troisième année de mécanique énergétique à l’Université de Toulouse et souhaitant poursuivre un master en modélisation et simulation numérique, je recherche une alternance en calcul mécanique au sein de votre équipe chez Capgemini Engineering.
	
	Passionné par la simulation physique des systèmes mécaniques, je dispose de bonnes bases en mécanique des solides, mécanique des fluides, transferts thermiques ainsi qu’en méthodes numériques et programmation Python. Je m’intéresse particulièrement aux activités d’ingénierie appliquée à la R\&D et à l’optimisation des structures mécaniques.
	
	Mon expérience au sein du projet Tls’e Racing m’a permis de réaliser des simulations par éléments finis avec ANSYS et Abaqus pour l’analyse structurale d’un conteneur de batterie de véhicule électrique soumis à des sollicitations mécaniques. J’ai également participé à la présentation des travaux de la division powertrain lors de la compétition Formula Student 2025 à Silverstone devant un jury international.
	
	Je souhaite approfondir mes compétences dans l’analyse des cahiers des charges, le dimensionnement analytique, la simulation par éléments finis ainsi que l’utilisation d’outils industriels tels que Patran-Nastran ou Hypermesh, en cohérence avec les missions proposées.
	
	Rejoindre Capgemini Engineering représente pour moi une opportunité de participer à des projets innovants de R\&D, d’évoluer dans un environnement technique stimulant et de contribuer à des solutions d’ingénierie à fort impact industriel et sociétal.
	
	Toujours motivé, rigoureux et passionné par la simulation numérique, je suis très enthousiaste à l’idée d’intégrer votre équipe d’ingénieurs en calcul mécanique pour développer à la fois mes compétences techniques et ma compréhension des problématiques industrielles.\\
	
	Je vous remercie pour l’attention portée à ma candidature et reste à votre disposition pour un entretien.
	
	Je vous prie d’agréer, Madame, Monsieur, l’expression de mes salutations distinguées. \\ 
	
	\vspace{1cm}
	
	\begin{flushright}
		\textbf{Friedly WOLI}
	\end{flushright}
	
\end{document}